\section{From Diagonal-Dominance Fingerprints to Lineage Verification}\label{sec:verification}

Trace concentration identifies valid trained pairings, not ancestry: independently trained models exhibit the same identity-aligned component. The verifier operates on the centered traceless remainder $E_\ell$, using trace concentration only to gate unreliable branches. The identity component is generic; $E_\ell$ is checkpoint-specific, surviving weight-preserving transformations but not reinitialization.

\subsection{Centered Residual Signatures and Block Alignment}\label{sec:pairing}

For each branch product $M_\ell$ we retain only the traceless remainder of decomposition~\eqref{eq:decomp}, normalized to a unit vector:
\begin{equation}\label{eq:centered-fingerprint}
R_\ell = M_\ell - \tfrac{\mathrm{tr}(M_\ell)}{d} \cdot I, \quad
\phi_\ell = \tfrac{\mathrm{vec}(R_\ell)}{\|\mathrm{vec}(R_\ell)\|_2}
\end{equation}
$R_\ell = E_\ell$ from~\eqref{eq:decomp}. The signature $\phi_\ell$ carries checkpoint-specific geometry that weight-preserving transformations inherit and independent training cannot reproduce. Because $M_\ell$ is invariant under hidden permutation and reciprocal rescaling, so is every signature built on it (Section~\ref{sec:exp-baselines}).

Given reference $A$ and suspect $B$ with $L$ blocks each, we form the similarity matrix $G_{ij} = \langle \phi^A_i, \phi^B_j \rangle$, gated by trace concentration (Appendix~\ref{app:verification}), and recover the block correspondence via the Hungarian algorithm~\citep{kuhn1955hungarian}:
\begin{equation}\label{eq:hungarian}
\pi^* = \arg\max_{\pi \in \mathcal{P}_L} \sum_{i} G_{i,\pi(i)}
\end{equation}
For derived checkpoints preserving block order, $\pi^*$ should be the identity; recovering it without metadata confirms that signatures track block identity. Complexity: branch products $O(Ld^2 h)$, similarity matrix $O(L^2 d^2)$, assignment $O(L^3)$.

\subsection{Lineage Score and Verification Test}\label{sec:lineage-score}

Diagonal dominance certifies trained residual structure, not ancestry: a verifier built on the dominance score alone achieves AUROC$=$0.417 (Appendix~\ref{app:cifar-resnet}, Figure~\ref{fig:lineage-roc}). To distinguish lineage, we compare the centered signatures $\phi_\ell$. The lineage score averages signature similarity over aligned branches:
\begin{equation}\label{eq:lineage-score}
\mathcal{L}(A, B) = \frac{1}{L} \sum_{\ell=1}^{L} \langle \phi^A_\ell, \phi^B_{\pi^*(\ell)} \rangle
\end{equation}
Each term is a cosine similarity, so $\mathcal{L} \in [-1,1]$ with $\mathcal{L}(A,A) = 1$. In practice, all observed scores are non-negative: related checkpoints score near 1, unrelated ones near 0, since independently initialized weights produce uncorrelated $E_\ell$. The score is symmetric.

We calibrate against a null distribution $\mathcal{N}_A = \{\mathcal{L}(A, U_j)\}$ from independently trained models $U_j$. The protocol returns \textsc{Related} if $\mathcal{L}(A,B)$ exceeds the maximum null score, and \textsc{Unrelated} otherwise:
\begin{equation}\label{eq:verdict}
V(A,B) = \begin{cases}
\textsc{Related} & \text{if } \mathcal{L}(A,B) > \max_{U_j} \mathcal{L}(A, U_j) \\
\textsc{Unrelated} & \text{otherwise}
\end{cases}
\end{equation}
With $n$ null models, this empirical threshold yields a conformal p-value of at most $1/(n+1)$; no distributional assumptions are required. We report AUROC as the primary metric, which sidesteps threshold selection entirely. Compatibility (same $L$ and $d$) is checked before scoring; the protocol returns \textsc{Incompatible} otherwise. Memory: $O(L^2 + d^2)$ beyond the checkpoints.
